\documentclass[11pt]{amsart}
\usepackage{amsmath,amssymb,amsthm}
\usepackage[margin=1.05in]{geometry}
\raggedbottom
\usepackage[hidelinks]{hyperref}
\usepackage{url}
\newtheorem{theorem}{Theorem}[section]
\newtheorem{proposition}[theorem]{Proposition}
\newtheorem{lemma}[theorem]{Lemma}
\newtheorem{corollary}[theorem]{Corollary}
\theoremstyle{remark}
\newtheorem{remark}[theorem]{Remark}
\newtheorem{conjecture}[theorem]{Conjecture}
\newcommand{\E}{\mathbb{E}}
\newcommand{\N}{\mathbb{N}}
\newcommand{\R}{\mathbb{R}}
\newcommand{\Prob}{\mathrm{Prob}}
\newcommand{\Lip}{L}

\title[The Wasserstein separation of the broken phases]{The Wasserstein separation of the broken phases:\\ a finite-dimensional reduction and the quadratic chaos bound}
\author{Ruqing Chen}
\address{GUT Geoservice Inc., Montreal, Canada}
\email{ruqing@hotmail.com}
\thanks{Paper XXIX of the series on localized Bost--Connes systems. The framework is that
of \cite{Main}, cited by DOI; Paper IV of the series, referred to by numeral, is an
unpublished companion, and the one fact used from it (the squeeze theorem) is re-proved in
place.}
\date{5 September 2026}

\begin{document}

\begin{abstract}
For a localized Bost--Connes system in its symmetry-broken phase, the Connes--Rieffel
distance $W_1(\varphi_+,\varphi_-)$ between the two extreme $\mathrm{KMS}_\beta$ states was
squeezed in Paper IV of the series between $w_0\Pi_\infty(\beta)$ and $w_0$, where $w_0$ is
the weight of the group direction and $\Pi_\infty(\beta)=\prod_j\frac{1-t_j}{1+t_j}$,
$t_j=\ell_j^{-\beta}$, is the Kakutani product of the order parameter. We sharpen both
sides. First, a finite-dimensional reduction: with $W_1^{(N)}$ the distance of the
$N$-coordinate truncation, an exactly computable second-order cone program,
\[
\Pi_{(N,\infty)}\,W_1^{(N)}\ \le\ W_1(\varphi_+,\varphi_-)\ \le\
\Pi_{(N,\infty)}\,W_1^{(N)}+2\rho_N,
\]
the lower bound exact on truncations and the tail variance
$\rho_N^2=2\sum_{j>N}w_j^2t_j/(1-t_j)^2$ explicit. Second,
testing against the degree-one odd Walsh chaos gives the closed-form lower bound
\[
W_1(\varphi_+,\varphi_-)\ \ge\
\Pi_\infty(\beta)\Bigl(w_0^2+4\sum_{j}\frac{w_j^2\,t_j^2}{(1-t_j)^2}\Bigr)^{1/2},
\]
valid whenever $w_0^2\ge2\sum_jw_j^2t_j/(1-t_j)$ and strictly sharper than
$w_0\Pi_\infty(\beta)$ at every finite $\beta$. The bound collapses onto the squeeze at the
deep-freeze limit $\beta\to\infty$ and provides an explicit envelope for the second-order
vanishing at $\beta_c$. We show that near-optimal test functions can be taken asymptotically
finite-level and odd under the symmetry, and we record the two questions that remain: a
matching upper bound, and whether the optimal profile is parity-measurable.
\end{abstract}

\maketitle

\section{The setting, and the squeeze re-proved}\label{sec:setting}

Throughout, $\beta$ lies in the symmetry-broken phase: $J$ is a countable index set of
primes $\ell_j$ with $t_j=\ell_j^{-\beta}$ and $\sum_jt_j<\infty$. The base space is
$X=G\times\prod_{j\in J}\overline\N$, with $G$ a compact abelian group carrying a fixed
character $\chi$ of order $2$ and a distinguished element $\sigma\in G$ with
$\chi(\sigma)=-1$. The symmetric measure is
$\nu_{\mathrm{sym}}=\mathrm{Haar}\otimes\bigotimes_j\mu_{t_j}$, where $\mu_t$ is the
geometric law $\mu_t(k)=(1-t)t^k$ on $\N$. Write $\varepsilon_j(x)=(-1)^{x_j}$, so
\[
m_j:=\E[\varepsilon_j]=\frac{1-t_j}{1+t_j}\in(0,1),
\qquad
\Pi_{(N,\infty)}:=\prod_{j>N}m_j,\qquad \Pi_\infty:=\Pi_{(0,\infty)} .
\]
Since $\sum_j\mu_{t_j}(x_j\ \mathrm{odd})=\sum_jt_j/(1+t_j)<\infty$, Borel--Cantelli makes
all but finitely many $\varepsilon_j$ equal to $1$ almost surely, so the order parameter
\[
\Psi(g,x)=\chi(g)\prod_{j\in J}\varepsilon_j(x_j)\in\{\pm1\}
\]
is defined $\nu_{\mathrm{sym}}$-a.e., and $\nu_{\mathrm{sym}}(\Psi=\pm1)=\tfrac12$ by the
$\chi$-symmetry of Haar measure. The extreme states are
$\varphi_\pm(f)=2\int f\,\mathbf1_{\{\Psi=\pm1\}}\,d\nu_{\mathrm{sym}}$, so that
\begin{equation}\label{eq:sep}
\varphi_+(f)-\varphi_-(f)=2\,\E[f\,\Psi].
\end{equation}
The Lipschitz seminorm, in the Rieffel formalism \cite{Rieffel}, is
\[
\Lip(f)=\Bigl\|\Bigl(w_0^{-2}|\Delta_0f|^2+\sum_{j\in J}w_j^{-2}|\Delta_jf|^2\Bigr)^{1/2}\Bigr\|_\infty,
\]
with $\Delta_0f(g,x)=f(g\sigma,x)-f(g,x)$ the group difference,
$\Delta_jf$ the one-step difference in the $j$-th coordinate, and weights satisfying
$\sum_jw_j^2<\infty$. The distance is
$W_1(\varphi_+,\varphi_-)=\sup\{\varphi_+(f)-\varphi_-(f):f\in C(X)_{sa},\ \Lip(f)\le1\}$.

\begin{lemma}[Squeeze, re-proved]\label{lem:squeeze}
$w_0\,\Pi_\infty(\beta)\ \le\ W_1(\varphi_+,\varphi_-)\ \le\ w_0$.
\end{lemma}

\begin{proof}
Upper: the map $T(g,x)=(g\sigma,x)$ preserves $\nu_{\mathrm{sym}}$ and flips $\Psi$, hence
transports $\nu_+:=2\mathbf1_{\{\Psi=1\}}\nu_{\mathrm{sym}}$ to $\nu_-$; for $\Lip(f)\le1$
one has $|f\circ T-f|=|\Delta_0f|\le w_0$ pointwise, so
$|\varphi_+(f)-\varphi_-(f)|=|\int(f-f\circ T)\,d\nu_+|\le w_0$. Lower: $f=\frac{w_0}{2}\chi$
has $\Delta_0f=-w_0\chi$, all other differences zero, so $\Lip(f)=1$, and \eqref{eq:sep}
gives $2\E[f\Psi]=w_0\,\E\bigl[\chi^2\textstyle\prod_j\varepsilon_j\bigr]=w_0\Pi_\infty$.
\end{proof}

\section{The finite-dimensional reduction}\label{sec:reduction}

Fix $N$ and let $\mathcal F_{\le N}$ be the $\sigma$-algebra generated by $g$ and
$x_1,\dots,x_N$, with $E_N$ the corresponding conditional expectation (integrating out
$x_j$, $j>N$). Write $\Psi_N=\chi\prod_{j\le N}\varepsilon_j$ and
$T_N=\prod_{j>N}\varepsilon_j$, so $\Psi=\Psi_NT_N$ with $T_N$ independent of
$\mathcal F_{\le N}$ and $\E[T_N]=\Pi_{(N,\infty)}$. Define the truncated distance
\[
W_1^{(N)}:=\sup\bigl\{2\,\E[f\,\Psi_N]:\ f\ \mathcal F_{\le N}\text{-measurable},\
\Lip(f)\le1\bigr\}.
\]

\begin{theorem}[Sandwich by truncations]\label{thm:reduction}
For every $N$,
\[
\Pi_{(N,\infty)}\,W_1^{(N)}\ \le\ W_1(\varphi_+,\varphi_-)\ \le\
\Pi_{(N,\infty)}\,W_1^{(N)}\ +\ 2\rho_N,
\quad\text{where }
\rho_N^2:=2\sum_{j>N}\frac{w_j^2\,t_j}{(1-t_j)^2}\to0 .
\]
The lower bound is exact on truncations: for an optimizer $f_N$ of $W_1^{(N)}$ the value
$2\,\E[f_N\Psi]$ equals $\Pi_{(N,\infty)}W_1^{(N)}$ exactly.
\end{theorem}

\begin{proof}
\emph{Lower.} For $\mathcal F_{\le N}$-measurable $f$, independence gives
$\E[f\Psi]=\E[f\Psi_N]\,\E[T_N]$, which is $\Pi_{(N,\infty)}\,\E[f\Psi_N]$; take the
supremum.

\emph{Upper.} Split $f=E_Nf+(f-E_Nf)$. For $j\le N$ the difference $\Delta_j$ commutes with
$E_N$, and $\Delta_j E_Nf=0$ for $j>N$; by the vector-valued Jensen inequality applied to
the $\ell^2$-norm inside the supremum, $\Lip(E_Nf)\le\Lip(f)\le1$, so the first term
contributes at most $\Pi_{(N,\infty)}W_1^{(N)}$ by the computation above. For the second
term, expand $f-E_Nf=\sum_{j>N}D_jf$ into martingale differences \cite[Ch.~3]{BLM} along the coordinates
$j>N$ (in any fixed order); the increments are orthogonal, and
$|D_jf|\le\E_{x_j'}\bigl|f(\dots,x_j,\dots)-f(\dots,x_j',\dots)\bigr|\le
w_j\,\E|x_j-x_j'|$ pointwise, since $\Lip(f)\le1$ forces $|\Delta_jf|\le w_j$ at every
point and the coordinate can be moved one step at a time. Hence
\[
\E|f-E_Nf|^2=\sum_{j>N}\E|D_jf|^2\le\sum_{j>N}w_j^2\,\E\bigl[(x_j-x_j')^2\bigr]
=\sum_{j>N}w_j^2\cdot2\,\mathrm{Var}(x_j)=\rho_N^2,
\]
using $\mathrm{Var}_{\mu_t}(x)=t/(1-t)^2$. Since $|\Psi|=1$, Cauchy--Schwarz bounds the
second term's contribution to \eqref{eq:sep} by $2\rho_N$.
\end{proof}

\begin{remark}[The truncation is a second-order cone program]\label{rem:socp}
$W_1^{(N)}$ maximizes the linear functional $2\E[f\Psi_N]$ of the finitely many values of
$f$ subject to one second-order cone constraint per point of the (truncated) configuration
space: after additionally truncating each coordinate to $x_j\le M$ --- an error controlled
by the geometric tails $\sum_{j\le N}w_jt_j^{M}/(1-t_j)$ through the same martingale
argument --- this is a standard SOCP with $O\bigl(2\cdot M^N\bigr)$ variables (the group
coordinate enters only through the two values of $\chi$), solvable exactly by interior-point
methods. Theorem~\ref{thm:reduction} converts any numerical solution into a certified
two-sided bound on $W_1$.
\end{remark}

\section{The quadratic chaos bound}\label{sec:quadratic}

We now evaluate the degree-one odd Walsh chaos exactly.

\begin{theorem}[Quadratic sum lower bound]\label{thm:quadratic}
Set $\kappa_j:=\dfrac{1-m_j}{m_j}=\dfrac{2t_j}{1-t_j}$, and suppose
\begin{equation}\label{eq:interior}
w_0^2\ \ge\ \sum_{j\in J}w_j^2\,\kappa_j .
\end{equation}
Then
\[
W_1(\varphi_+,\varphi_-)\ \ge\
\Pi_\infty(\beta)\Bigl(w_0^2+\sum_{j\in J}w_j^2\,\kappa_j^2\Bigr)^{1/2}
=\Pi_\infty(\beta)\Bigl(w_0^2+4\sum_{j\in J}\tfrac{w_j^2\,t_j^2}{(1-t_j)^2}\Bigr)^{1/2},
\]
the series converging by $\sum_jw_j^2<\infty$ and $t_j\to0$. The bound is strictly larger
than the squeeze lower bound $w_0\Pi_\infty(\beta)$ at every $\beta<\infty$, and collapses
onto it as $\beta\to\infty$.
\end{theorem}

\begin{proof}
Take the ansatz $f(g,x)=\chi(g)\bigl(c+\sum_jr_j\varepsilon_j(x_j)\bigr)$ with
$c,r_j\ge0$ and $\sum_jr_j<\infty$. Its differences are
$\Delta_0f=-2\chi(c+\sum_jr_j\varepsilon_j)$ and
$\Delta_jf=\chi\,r_j\,\Delta_j\varepsilon_j$ with $|\Delta_j\varepsilon_j|=2$, so the
Lipschitz constraint reads, pointwise in the signs $\varepsilon\in\{\pm1\}^J$,
\[
\frac{4\bigl(c+\sum_jr_j\varepsilon_j\bigr)^2}{w_0^2}+\sum_j\frac{4r_j^2}{w_j^2}\ \le\ 1 .
\]
The supremum over sign patterns is attained at $\varepsilon_j\equiv+1$, so with
$s:=c+\sum_jr_j$ the constraint is exactly the ellipsoid
$s^2/(w_0/2)^2+\sum_jr_j^2/(w_j/2)^2\le1$. Independence and $\varepsilon_j^2=1$ give the
objective
\[
2\,\E[f\Psi]=2\Bigl(c\,\Pi_\infty+\sum_jr_j\,\frac{\Pi_\infty}{m_j}\Bigr)
=2\Pi_\infty\Bigl[s+\sum_jr_j\kappa_j\Bigr],
\]
after substituting $c=s-\sum_jr_j$. Maximizing the linear functional
$s+\sum_jr_j\kappa_j$ over the ellipsoid gives the value
$\tfrac12\bigl(w_0^2+\sum_jw_j^2\kappa_j^2\bigr)^{1/2}$, attained at
$s^*=\tfrac{(w_0/2)^2}{R}$, $r_j^*=\tfrac{(w_j/2)^2\kappa_j}{R}$ with
$R=\tfrac12(w_0^2+\sum w_j^2\kappa_j^2)^{1/2}$; the sign conditions $r_j^*\ge0$ hold, and
$c^*=s^*-\sum_jr_j^*\ge0$ is exactly \eqref{eq:interior}. Multiplying by $2\Pi_\infty$
yields the bound. Strictness over $w_0\Pi_\infty$ holds whenever some $t_j>0$, i.e.\ at
every finite $\beta$; as $\beta\to\infty$ every $\kappa_j\to0$ and the correction vanishes.
\end{proof}

\begin{remark}[The boundary case]\label{rem:boundary}
If \eqref{eq:interior} fails --- possible near $\beta_c$ when the weights decay slowly ---
the maximum of the same linear functional is taken on the face $c=0$ of the cone, i.e.\ on
the ellipsoid intersected with $\{\sum_jr_j=s\}$; this is again a closed-form Lagrange
computation, and the resulting bound still dominates $w_0\Pi_\infty(\beta)$, but we do not
display the longer formula. In either regime the theorem provides an explicit lower
envelope for the second-order vanishing of $W_1$ at $\beta_c$ established in Paper IV.
\end{remark}

\section{Structure of optimizers, and what remains}\label{sec:structure}

\begin{proposition}[Odd, asymptotically finite-level]\label{prop:structure}
The supremum defining $W_1$ is unchanged if restricted to functions that are odd under
$T(g,x)=(g\sigma,x)$, and for every $\epsilon>0$ it is attained within $\epsilon$ by an
$\mathcal F_{\le N}$-measurable function as soon as $2\rho_N\le\epsilon$.
\end{proposition}

\begin{proof}
The odd part $f^{\mathrm o}=\tfrac12(f-f\circ T)$ satisfies
$\E[f^{\mathrm o}\Psi]=\E[f\Psi]$ (since $\Psi\circ T=-\Psi$ and $\nu_{\mathrm{sym}}$ is
$T$-invariant) and $\Lip(f^{\mathrm o})\le\Lip(f)$ (each difference of $f^{\mathrm o}$ is
the average of two differences of $f$). The second statement is
Theorem~\ref{thm:reduction} with $E_Nf$ in place of $f$.
\end{proof}

\begin{conjecture}\label{conj:hierarchy}
The increasing sequence of values of the odd Walsh-chaos programs --- degree one being
Theorem~\ref{thm:quadratic}, degree $d$ allowing coefficients
$r_A\prod_{j\in A}\varepsilon_j$ with $|A|\le d$ --- converges to $W_1(\varphi_+,\varphi_-)$.
\end{conjecture}

\begin{remark}[Numerical verdict]\label{rem:numerics}
The programs of Remark~\ref{rem:socp} were solved for a toy model ($w_0=1$,
$w_j=(j+1)^{-1}$, $t_j=p_j^{-1}$, $N\le3$; the feasible set is convex and the objective
linear, so any KKT point is the global maximum, and constraint violations were at machine
precision). Two findings, scripts and outputs in \texttt{code/conj42/} of the repository.
First, the degree-one bound of Theorem~\ref{thm:quadratic} is far from sharp: the exact
$2^N$-point parity program --- whose value is the limit of the odd Walsh-chaos hierarchy,
since a forward step always flips parity and the constraint closes on parity classes
without any truncation --- exceeds it by $26\%$ at $N=2$ and $39\%$ at $N=3$. Higher
chaos participates massively. Second, and against Conjecture~\ref{conj:hierarchy}: the
full grid optimum exceeds the parity optimum, in matched truncations, by a stable
$2$--$4\times10^{-3}$ that is independent of the grid depth $M$ at $N\ge2$, while at
$N=1$ (where parity-measurable and degree-one coincide) the same gap vanishes to
$10^{-15}$ once the boundary parity is matched --- a control showing the $N\ge2$ gap is a
bulk effect, not a boundary artifact. The evidence is therefore that the optimal profile
is \emph{not} parity-measurable and that the odd-chaos hierarchy converges to a value
strictly below $W_1^{(N)}$: Conjecture~\ref{conj:hierarchy} is expected to be false as
stated, and the parity value should be read as a computable, certified lower bound rather
than the answer.
\end{remark}

Two questions remain, and they are the genuine content left in this problem. First, a
matching \emph{upper} bound: the transport map of Lemma~\ref{lem:squeeze} moves mass only
in the group direction; any improvement must couple the group move to the parity
configuration, an $x$-adaptive coupling whose cost functional is not yet controlled.
Second, \emph{parity measurability} (on which Remark~\ref{rem:numerics} already gives a numerical answer): the order parameter sees only the parities
$\varepsilon_j$, but conditioning a test function on the parity $\sigma$-algebra does not
obviously contract $\Lip$, because the one-step differences do not respect the parity
fibres; whether genuinely non-parity-measurable test functions can strictly increase the
supremum is open. A numerical test of Conjecture~\ref{conj:hierarchy} is available today:
Remark~\ref{rem:socp} certifies two-sided bounds for small $N$, to be compared with the
closed form of Theorem~\ref{thm:quadratic}.

\begin{thebibliography}{9}
\bibitem{Main} R. Chen, \emph{KMS state uniqueness and phase transitions in localized Bost--Connes systems}, preprint, 2026, \newline\url{https://doi.org/10.5281/zenodo.22152101}.
\bibitem{Rieffel} M. A. Rieffel, \emph{Metrics on state spaces}, Doc. Math. 4 (1999), 559--600.
\bibitem{BLM} S. Boucheron, G. Lugosi, P. Massart, \emph{Concentration Inequalities: A Nonasymptotic Theory of Independence}, Oxford University Press, 2013.
\end{thebibliography}

\end{document}
